\chapter{Arbitrage and risk-neutral pricing}
\label{ch:arbitrage}

\begin{question}
A stock will be worth one of two prices tomorrow. What is the fair price today of an option on it, and why can no other price survive?
\end{question}

\begin{whatfor}
\begin{itemize}
  \item Discounting: money later is worth less than money now.
  \item Prices an option in a one-period, two-price market by building a riskless holding, and finds the only price that allows no sure profit.
  \item Gets the same price from the risk-neutral probability, under which every bet has expected return $0$.
  \item States the arbitrage theorem for any finite experiment.
\end{itemize}
\skipline{fs:arbitrage}
\end{whatfor}

\begin{prereq}
From Part~0: probability on a finite set of outcomes, expectation. From Chapter~5: the binomial tree and the martingale probability $p=\frac{1-d}{u-d}$. (This chapter needs nothing from Parts~2--3.)
\end{prereq}

%=====================================================================
\section{Money now and money later}

An amount $v$ received at time $t$ in the future is not worth as much as $v$ received now: money now can earn interest. With a constant rate $\alpha$, the amount $v$ at time $t$ is worth $v\,e^{-\alpha t}$ today. The number $\alpha$ is the \dfn{discount rate} and $e^{-\alpha t}$ the \dfn{discount factor}. For example Rs~100 a year from now, at $\alpha=0.1$, is worth $100e^{-0.1}\approx\text{Rs }90.48$ today.

In this chapter all prices are quoted in present value, so no further discounting is needed.

%=====================================================================
\section{One stock, one option, two prices}

\begin{setup}[The one-period market]
A stock costs $100$ today (time $0$). At time $1$ its price is either $200$ or $50$. An \dfn{option} costs $C$ today and gives the right, but not the obligation, to buy one share at time $1$ for the \dfn{strike price} $K=150$.
\end{setup}

At time $1$ the option is worth $200-150=50$ if the price is $200$ (use it), and $0$ if the price is $50$ (buying at $150$ a share worth $50$ would be a loss, so do not use it).

\begin{pic}
\begin{tikzpicture}[x=3cm,y=1.1cm,every node/.style={font=\small}]
  \node[draw,rounded corners,fill=defc!8] (s0) at (0,0) {stock $100$};
  \node[draw,rounded corners,fill=thmc!10] (su) at (1,1) {stock $200$};
  \node[draw,rounded corners,fill=trapc!8] (sd) at (1,-1) {stock $50$};
  \draw[-{Latex}] (s0)--(su); \draw[-{Latex}] (s0)--(sd);
  \node[anchor=west] at (1.35,1) {option pays $200-150=50$};
  \node[anchor=west] at (1.35,-1) {option pays $0$};
  \node[prec] at (0,-1.8) {$t=0$}; \node[prec] at (1,-1.8) {$t=1$};
  \node[anchor=east] at (-0.3,0) {option costs $C$};
\end{tikzpicture}
\picnote{The whole market: one stock with two possible prices tomorrow, and one option whose payoff is fixed by the stock price.}
\end{pic}

\paragraph{A holding with no risk.} Buy $x$ shares and $y$ options at time $0$ (a negative number means selling). At time $1$ the holding is worth
\[
  \text{Value}=\begin{cases}200x+50y,&\text{price }200,\\ 50x,&\text{price }50.\end{cases}
\]
Choose $y$ so the two values are equal: $200x+50y=50x$ gives $y=-3x$. Then the holding is worth $50x$ whatever happens. It cost $100x+Cy=100x-3xC$, so
\[
  \text{gain}=50x-(100x-3xC)=x\,(3C-50).
\]
If $3C>50$, take $x>0$; if $3C<50$, take $x<0$. Either way the gain is positive \emph{for sure}.

\begin{example}[Two option prices, two sure profits]
\label{ex:arb}
One overpriced and one underpriced option.
\begin{itemize}
  \item $C=20$: buy $1$ share, sell $3$ options ($x=1$, $y=-3$). Cost $100-60=40$; value at time $1$ is $50$ for sure. Profit Rs~$10$, with no risk.
  \item $C=15$: sell $1$ share, buy $3$ options ($x=-1$, $y=3$). You receive $100-45=55$; at time $1$ you owe $50$. Profit Rs~$5$, with no risk.
\end{itemize}
\end{example}

\begin{definition}[Arbitrage]
An \dfn{arbitrage} is a way of buying and selling that is sure to make a profit: a risk-free profit.
\end{definition}

The only option price that allows no arbitrage is $3C=50$, i.e.\ $C=\frac{50}{3}\approx16.67$.

\begin{pic}
\begin{tikzpicture}
\begin{axis}[width=9cm,height=4.4cm,xmin=10,xmax=25,ymin=-12,ymax=12,xlabel={option price $C$},ylabel={sure gain},domain=10:25,legend style={font=\scriptsize,at={(0.02,0.98)},anchor=north west}]
  \addplot[defc,very thick] {3*x-50}; \addlegendentry{$x=1$: buy share, sell 3 options}
  \addplot[trapc,very thick] {50-3*x}; \addlegendentry{$x=-1$: sell share, buy 3 options}
  \addplot[only marks,mark=*] coordinates {(50/3,0)};
  \node[font=\scriptsize,anchor=north] at (axis cs:50/3,-1) {$C=\frac{50}3$};
  \addplot[prec,dotted] {0};
\end{axis}
\end{tikzpicture}
\picnote{At every option price except $\frac{50}{3}$ one of the two schemes makes a sure profit. Only at $C=\frac{50}{3}$ do both gains vanish.}
\end{pic}

%=====================================================================
\section{The risk-neutral probability}

Look at the same market through probabilities. Suppose the price goes to $200$ with probability $p$ and to $50$ with probability $1-p$. Ask for a $p$ under which each bet has expected return $0$.
\begin{itemize}
  \item Stock: the return on one share is $100$ or $-50$. $100p-50(1-p)=0$, i.e.\ $100=200p+50(1-p)$, gives $p=\frac13$. Today's price equals the expected price tomorrow.
  \item Option: the return is $50-C$ or $-C$. $(50-C)\frac13-C\cdot\frac23=0$ gives $C=50\cdot\frac13=\frac{50}{3}$.
\end{itemize}
The same price as the no-arbitrage argument.

\begin{definition}[Risk-neutral probability]
A probability on the outcomes under which every available bet has expected return $0$ (equivalently, every price today is the expected value of its payoff) is a \dfn{risk-neutral probability}.
\end{definition}

So: \emph{the no-arbitrage price of the option is the expected payoff under the risk-neutral probability.} The real-world probability of an up-move never entered. With $p^*=\frac13$ here, $C=\E_{P^*}(\text{payoff})=\frac13\cdot50$.

\paragraph{Three possible prices.} Let the stock go from $100$ to $300$, $200$ or $50$ with probabilities $p_1,p_2,1-p_1-p_2$ (all positive), and keep $K=150$, so the option pays $150$, $50$, $0$. The two conditions become
\[
  100=300p_1+200p_2+50(1-p_1-p_2)\quad(\text{i.e. }5p_1+3p_2=1),\qquad C=150p_1+50p_2 .
\]
One equation, two unknowns: many risk-neutral probabilities. Substituting $p_2=\frac{1-5p_1}{3}$ gives $C=\frac{50}{3}+\frac{200}{3}p_1$ with $0<p_1<\frac15$, so every $C$ in $\big(\frac{50}3,\,30\big)$ is consistent with no arbitrage. With three outcomes and only two assets, the option is not pinned down to one price. (Chapter~20 explains this: the market is not complete.)

%=====================================================================
\section{The arbitrage theorem}

\begin{setup}[A general betting experiment]
An experiment has outcomes $j\in\{1,\dots,m\}$. There are $n$ wagers; a unit bet on wager $i$ returns $r_i(j)$ if the outcome is $j$. A \dfn{betting scheme} $x=(x_1,\dots,x_n)$ bets $x_i$ on wager $i$ and returns $\sum_ix_ir_i(j)$ when the outcome is $j$.
\end{setup}

Think of the returns as a table: row $i$ is a wager, column $j$ an outcome. For the two-price market with $C=20$: the stock row is $(100,-50)$ and the option row is $(30,-20)$.

\begin{theorem}[Arbitrage theorem]
\label{thm:arbitrage}
Exactly one of the following holds:
\begin{enumerate}[label=(\roman*)]
  \item (\emph{no arbitrage}) there is a probability vector $p=(p_1,\dots,p_m)$, $p_j\ge0$, $\sum_jp_j=1$, with $\displaystyle\sum_{j=1}^mp_j\,r_i(j)=0$ for every wager $i$;
  \item (\emph{arbitrage}) there is a betting scheme $x$ with $\displaystyle\sum_{i=1}^nx_i\,r_i(j)>0$ for every outcome $j$.
\end{enumerate}
\end{theorem}
\noproof

Read (i) across each row (every wager has expected return $0$ under $p$) and (ii) down each column (every outcome gives a positive total). Both cannot hold: if they did, averaging the column totals of (ii) with the weights $p_j$ would give a positive number, but by (i) it equals $\sum_ix_i\sum_jp_jr_i(j)=0$.

\begin{example}[Checking the theorem]
With $C=20$: $x=(1,-3)$ gives $10$ in both columns, so (ii) holds. And (i) fails: the stock row forces $p=(\frac13,\frac23)$, but then the option row gives $30\cdot\frac13-20\cdot\frac23=-\frac{10}3\ne0$. With $C=\frac{50}3$, $p=(\frac13,\frac23)$ satisfies (i), so no sure-win scheme exists.
\end{example}

%=====================================================================
\formulasheet{fs:arbitrage}

\begin{mastercard}[label=card:oneperiod]{No-arbitrage price in a one-period market}
A stock costs $S_0$ today and is worth $S_j$ tomorrow in state $j=1,\dots,m$; a claim pays $X_j$ in state $j$; $D$ is the discount factor ($D=1$ if prices are in present value, $D=\frac1{1+r}$ with interest $r$ per period). Find the risk-neutral probabilities $p^*_j>0$, $\sum_jp^*_j=1$, from the stock, then price the claim:
\[
  \boxed{\begin{gathered}S_0=D\sum_jp^*_jS_j\quad\Longrightarrow\quad\text{price}=D\sum_jp^*_jX_j;\\
  m=2\ (S_u>S_d):\quad p^*=\frac{S_0/D-S_d}{S_u-S_d};\\ \text{hedge: }\Delta=\frac{X_u-X_d}{S_u-S_d}\ \text{shares and a bond worth }D(X_u-\Delta S_u)\text{ today}.\end{gathered}}
\]
Any other price allows an arbitrage (a sure profit); with more states than assets the $p^*$ are not unique and the price is only pinned to an interval.
\begin{plugrows}
  \plugrow{Two prices, present values}{$D=1$; $X=(S-K)^+$}{$p^*=\frac{S_0-S_d}{S_u-S_d}$, price $p^*(S_u-K)^+$}{$S_d<S_0<S_u$}{$100\to200/50$, $K=150$: $p^*=\frac13$, $C=\frac{50}3$; $75\to200/50$, $K=120$: $\frac{40}3$; $50\to100/25$, $K=85$: $5$}
  \plugrow{With interest}{$D=\frac1{1+r}$, $S_u=uS_0$, $S_d=dS_0$}{$p^*=\frac{1+r-d}{u-d}$}{$d<1+r<u$}{$S_0=100$, $u=1.2$, $d=0.8$, $r=0.05$, $K=100$: $p^*=0.625$, $C=\frac{0.625\cdot20}{1.05}\approx11.90$}
  \plugrow{The replicating holding}{second line}{$\Delta S_0+D(X_u-\Delta S_u)$ equals the price}{two states}{same: $\Delta=0.5$, bond $-38.10$, $50-38.10=11.90$}
  \plugrow{Building the arbitrage}{if the market price $C\ne$ fair price $C^*$: sell the dear one, buy the cheap one in the hedge ratio}{sure gain $=$ (price difference) per claim, at time 0 in present value}{two states}{$100\to200/50$, $K=150$, $C=20$: sell 3 options, buy 1 share: gain $3(20-\frac{50}3)=10$}
  \plugrow{More states than assets}{solve the stock equation with one free $p^*$; let it range over values keeping every $p^*_j>0$}{an interval of no-arbitrage prices}{all $p^*_j\in(0,1)$}{$100\to300/200/50$, $K=150$: $C=\frac{50}3+\frac{200}3p_1$, $p_1\in(0,\frac15)$: $C\in(16.67,\,30)$}
  \plugrow{The arbitrage theorem}{returns $r_i(j)$ of wager $i$ in outcome $j$}{either some $p$ gives every wager expected return $0$, or some betting scheme wins in every outcome; never both}{finitely many outcomes}{$C=20$: scheme $(1,-3)$ wins $10$ in both states}
\end{plugrows}
\end{mastercard}
